% Draft -- not for distribution.
% Source of truth for every statement: notes/dcltheorem_results.md (Theorem D', with the
% r13 fixes) and notes/multimod_results.md (Theorems A, A', B, C).
\documentclass[11pt]{amsart}

\usepackage[utf8]{inputenc}
\usepackage[T1]{fontenc}
\usepackage{amsmath,amssymb,amsthm}
\usepackage{mathtools}
\usepackage{booktabs}
\usepackage{enumitem}
\usepackage[margin=1.1in]{geometry}
\usepackage{xcolor}
\usepackage[colorlinks=true,linkcolor=blue!50!black,citecolor=green!35!black,urlcolor=blue!50!black]{hyperref}

\newtheorem{theorem}{Theorem}[section]
\newtheorem{lemma}[theorem]{Lemma}
\newtheorem{proposition}[theorem]{Proposition}
\newtheorem{corollary}[theorem]{Corollary}
\newtheorem{conjecture}[theorem]{Conjecture}
\theoremstyle{definition}
\newtheorem{definition}[theorem]{Definition}
\newtheorem{remark}[theorem]{Remark}

\newcommand{\Z}{\mathbb{Z}}
\newcommand{\N}{\mathbb{N}}
\newcommand{\R}{\mathbb{R}}
\newcommand{\Sn}{\mathcal{S}_n}
\newcommand{\Pk}{\mathcal{P}_k}
\newcommand{\ord}{\operatorname{ord}}
\newcommand{\mbar}{\bar m_k}
\newcommand{\spanbar}{\overline{\mathrm{span}}}
\newcommand{\DCL}{\textsf{DCL-det}}
\newcommand{\polylog}{\operatorname{polylog}}

\title[Deterministic factoring of almost all balanced semiprimes]{A deterministic
algorithm that factors almost all balanced semiprimes\\ in time $N^{1/6}(\log N)^{O(1)}$}

\author{[Author]}
\address{[Address]}
\email{[Email]}

\date{Draft -- not for distribution. Version of 18 October 2026.}

\begin{document}

\begin{abstract}
We give a deterministic algorithm that receives only an $n$-bit integer $N$, uses
$N^{1/6+5\varepsilon/6}(\log N)^{O(1)}$ bit operations on every input, never outputs anything
other than a proper divisor of $N$ or \texttt{FAIL}, and factors all but an
$O(N^{-\varepsilon})$ fraction of the semiprimes $N=pq$ with primes $p<q\le 2p$
(for any fixed $\varepsilon\in(0,1/60]$ and $n\ge n_0(\varepsilon)$). Taking polylogarithmic
weight heights instead gives cost $N^{1/6}(\log N)^{O(1)}$ with an exceptional fraction
$O((\log N)^{-A})$ for any fixed $A$ (the polylogarithmic exponent depends on $A$). The
best proved deterministic exponent for arbitrary inputs remains $1/5$ (Harvey;
Harvey--Hittmeir). The algorithm is a Lehman-type weighted collision search. It takes the
McKee--Pinch divisibility event $\ell\mid Bp-Aq$ over all primes $\ell\equiv3\pmod4$ in four
dyadic bands near $N^{1/6}$ and all small coprime weights $(A,B)$, together with the
Harvey--Hittmeir modulus-squared lock. Within each band, one baby-step list is shared by all
band primes through cofactor-powered handles. Three ingredients carry the proof. The
collision modulo $p$ is an exact identity, so no lower bound on any multiplicative order is
needed. Four disjoint bands guarantee one band coprime to $q-1$, and on that band a
two-sided (useless) collision forces a rare congruence modulo $\ord_q(2)$. The arithmetic
large sieve bounds the uncovered pairs. For context we record a lower bound of order
$N^{1/6}$ in a one-base sumset collision-search model; it applies only to families that cover
\emph{every} integer candidate for a fixed $N$, and so does not bound average-case algorithms
such as ours. We also report pre-registered experiments. Over 48--80 bits the measured cost
grows more slowly than that of a calibrated Harvey--Hittmeir implementation (slope $0.203$
against $0.226$ in $\log_2$ cost per bit), but at 40--80 bits Harvey--Hittmeir is still
$2.3$--$4.9$ times cheaper. The result says nothing about the worst case. Every ingredient is
prior art; we found no prior formulation of this combination.
\end{abstract}

\maketitle

\begin{center}
\fbox{\parbox{0.9\textwidth}{\small\textbf{Draft --- not for distribution.} This draft was
written with substantial AI assistance (see the Acknowledgements). Its proofs have been
checked by AI referees and have not yet been checked by a human referee. Status labels
(\emph{proved}, \emph{experimentally supported}, \emph{conjectured}, UNVERIFIED) are used
throughout.}}
\end{center}

\tableofcontents

%======================================================================
\section{Introduction and statement of results}
%======================================================================

\subsection{Background}
Deterministic integer factoring with rigorous bounds has a short list of exponents. Trial
division costs $N^{1/2}$. Lehman's method \cite{lehman1974} and the Pollard--Strassen
method \cite{pollard1974,strassen1976} give $N^{1/3+o(1)}$ and $N^{1/4+o(1)}$, and the
latter was refined by Bostan--Gaudry--Schost \cite{bgs2007} and Costa--Harvey
\cite{costaharvey}. Hittmeir \cite{hittmeir2018} saved a subexponential factor at exponent
$1/4$, and then proved exponent $2/9$ \cite{hittmeir2021} with a time--space trade-off for
Lehman's method. Harvey \cite{harvey2021} proved exponent $1/5$, and Harvey--Hittmeir
\cite{hh2022} improved the logarithmic factor; Gao, Feng, Hu and Pan \cite{gao2025} improve
logarithmic factors again for balanced semiprimes, at the same exponent $1/5$. Below $1/5$,
the known routes are conditional: Umans and Wang \cite{umanswang2025} showed that a
number-theoretic divisor-cover conjecture would give exponent $1/6$, and He and Sahai
\cite{hesahai2026} refuted that conjecture's one-dimensional arithmetic-progression form.
Carella \cite{carella} claims a deterministic $O(N^{1/6+\varepsilon})$ algorithm; we have not
verified that claim (Section~\ref{sec:prior}).

This paper is about the \emph{average} case over the hardest natural population, semiprimes
with two factors of comparable size. For unrestricted semiprimes the average case is easy:
almost all of them have a factor below $N^{\delta}$, so trial division handles them (see
\cite{crisanerban} for the counting). Balanced semiprimes are the opposite. Every method
above is analysed in the worst case, and we know of no proved average-case exponent below
$1/5$ for them.

\subsection{Population and main result}
Let $n$ be the bit length, $X=2^n$, $P=2^{n/2}$, and
\[
  \Sn=\{(p,q): p<q\le 2p,\ p,q\ \text{prime},\ 2^{n-1}\le pq<2^n\}.
\]
We count $\Sn$ by pairs. By unique factorisation this is the same as counting the
semiprimes $N=pq$. For $(p,q)\in\Sn$ we have $p\in[\sqrt{N/2},\sqrt N]\subset[P/2,P]$ and
$q\in(P/\sqrt2,\sqrt2P)$. The prime number theorem gives $|\Sn|\ge c_0P^2/n^2$ for an
absolute $c_0>0$ and all large $n$.

\begin{theorem}[Theorem D$'$; proved]\label{thm:main}
Fix $\varepsilon\in(0,1/60]$. There are $n_0(\varepsilon)$ and absolute constants $C,c$
such that for every $n\ge n_0$ the deterministic algorithm $\DCL_{n,\varepsilon}$ of
Section~\ref{sec:alg} has the following properties.
\begin{enumerate}[label=(\roman*)]
\item \emph{Correctness.} On every input $N$ it outputs either a proper divisor of $N$
  (an integer strictly between $1$ and $N$ that divides $N$) or \texttt{FAIL}.
\item \emph{Cost.} On every $n$-bit input it uses $O\bigl(n^{c}\,2^{(1/6+5\varepsilon/6)n}\bigr)
  = N^{1/6+5\varepsilon/6}(\log N)^{O(1)}$ bit operations.
\item \emph{Success.} It outputs \texttt{FAIL} on at most $C\,2^{-\varepsilon n}|\Sn|$ pairs
  $(p,q)\in\Sn$.
\end{enumerate}
In particular, for every $\delta\in(0,1/72]$ there is a deterministic algorithm of cost
$N^{1/6+\delta+o(1)}$ whose exceptional fraction on $\Sn$ is $O(N^{-6\delta/5})$.
\end{theorem}

\begin{corollary}[Corollary D$''$; proved]\label{cor:polylog}
For every fixed $a>0$, the same algorithm with weight height $H=\lceil n^{1+a}\rceil$ costs
$N^{1/6}(\log N)^{O(1+a)}$ bit operations and fails on $O((\log N)^{-2a})$ of $\Sn$.
\end{corollary}

\begin{remark}[Scope]\label{rem:scope}
We state the scope plainly, because each item is easy to over-read.
\begin{enumerate}[label=(\alph*)]
\item The statement is about the true factor of an input drawn uniformly from $\Sn$, with
  exceptions counted by pairs. It says nothing about any particular $N$. In particular it
  is \emph{not} a worst-case bound, and it does not show that, for almost every $N$, every
  integer candidate $p$ is covered (the ``full-candidate'' covering model of
  Section~\ref{sec:lower}).
\item Only ratios $q/p\le 2$ are covered. Ratios $q/p>2$ (and unbalanced semiprimes in
  general) are not.
\item ``Never wrong'' means: every output is validated as a proper divisor. The algorithm may
  output \texttt{FAIL} on inputs outside $\Sn$ and on the exceptional pairs.
\item The theorem is asymptotic ($n\ge n_0(\varepsilon)$). At practical sizes the algorithm
  is slower than Harvey--Hittmeir (Section~\ref{sec:experiments}).
\end{enumerate}
\end{remark}

\subsection{How it works, in one paragraph}
Fix a prime $\ell$ and coprime weights $(A,B)$. If $\ell\mid Bp-Aq$, then $p\bmod\ell$ is
one of the two square roots of $AN/B$ modulo $\ell$, and these roots are computable from
$N$. On that residue class the weighted sum $Bp+Aq$ is determined modulo $\ell^2$ (the
Harvey--Hittmeir lock), so $g=(Bp+Aq-c)/\ell^2$ is an integer in a window of length about
$H\sqrt N/\ell^2$. Exponentiating, $2^{(L/\ell^2)(B+AN-c)}$ equals $\gamma^{g}$ modulo $p$,
with $\gamma=2^{L}$ and $L=\prod\ell^2$ over the band. Every target, for every prime of the
band, is therefore one giant step against the band's single baby list
$\gamma^0,\dots,\gamma^{m-1}$, and a
product-tree/remainder-tree evaluation detects collisions modulo the unknown $p$. Primes
near $N^{1/6}$ balance the baby list ($\approx H\sqrt N/\ell^2$) against the number of
targets ($\approx H^2\ell$). Coverage is a sieve problem in $p$ for fixed $q$, and the large
sieve shows that the true $p$ is caught for almost every pair. The collision modulo $p$ is exact, so no order
lemma is needed there. It only remains to show that the true giant does not \emph{also}
collide modulo $q$. Using four disjoint bands, one of them is coprime to $q-1$. On that band
a two-sided collision forces $\ord_q(2)$ to divide a specific small expression, which a CRT
count shows to be rare.

\subsection{A lower bound in the collision-search model}
For context, Section~\ref{sec:lower} records a lower bound in a model whose targets include
those of Lehman, Hittmeir, Harvey, Harvey--Hittmeir and the present algorithm (one base,
polynomial exponent maps of bounded degree, unit residue-class locks modulo any number of
moduli, sumset baby-step/giant-step collision detection). In that model every family that
covers \emph{all} integer candidates in the balanced range, for a fixed $N$, costs at least
$cN^{1/6}$ (Theorem~\ref{thm:A}). With $N^{o(1)}$ distinct moduli, or with targets of degree
at least $2$, the bound is $N^{1/5-o(1)}$ (Theorems~\ref{thm:Aprime} and~\ref{thm:B}). This is
a floor for full-candidate covering in this model only. It is not a lower bound for
factoring, and it does not bound average-case algorithms, which (like ours) need only reach
the true factor for most inputs.

\subsection{What is new}
Each ingredient is known (Section~\ref{sec:prior}). The divisibility event
$\lambda p\equiv\pm\mu q\pmod k$ with small multipliers, and a heuristic balanced method
over primes $k\approx N^{1/4}$, are in McKee and Pinch \cite[\S5]{mckeepinch1996}. The
modulus-squared lock is in McKee \cite{mckee1999} and Harvey--Hittmeir
\cite[Prop.~3.3]{hh2022}. Shared collision search with multipoint evaluation is in Harvey
\cite[Alg.~4.1--4.2]{harvey2021}. The large sieve is Montgomery's
\cite{montgomery1968,mv1973}. What this paper adds is the following combination:
\begin{enumerate}[label=(\roman*)]
\item a band of primes at $N^{1/6}H^{-1/3}$ instead of $N^{1/4}$,
\item made possible by one baby list per band, shared across all primes of the band via
  cofactor-powered handles, so that each target costs one giant step;
\item a proof (rather than a probability heuristic) that the event occurs for all but
  $O(N^{-\varepsilon})$ of the pairs: the prime-pair large sieve with growing weight height;
\item one-sided extraction with no order hypothesis: the clean-band argument, the
  elementary bound on primes with small $\ord_q(2)$, and a CRT count of two-sided true
  collisions;
\item a full bit-complexity accounting; for context, a model lower bound for full-candidate
  covering (which does not apply to average-case algorithms).
\end{enumerate}
We found no prior formulation of this combination.

\subsection{Organisation}
Section~\ref{sec:prior} reviews prior art. Section~\ref{sec:alg} defines the algorithm.
Section~\ref{sec:proofs} proves Theorem~\ref{thm:main} and Corollary~\ref{cor:polylog}.
Section~\ref{sec:lower} gives the lower bound in the collision-search model.
Section~\ref{sec:experiments} reports pre-registered experiments. Section~\ref{sec:open}
lists limitations and open problems.

%======================================================================
\section{Prior art}\label{sec:prior}
%======================================================================

We separate what we read ourselves from what we know only through AI literature sweeps.
We read McKee--Pinch \S5 ourselves, from the scans. The other descriptions below come from
the GPT-6.1-sol sweeps described in the Acknowledgements, which checked the statements
against the sources; they are marked \emph{[sweep]} and were not re-read by us. Where even
the sweep did not check a source, we say UNVERIFIED.

\paragraph{Lehman and weighted Fermat searches.}
Lehman \cite{lehman1974} searches for $4kN=x^2-y^2$ over small multipliers $k$ and proves a
deterministic $N^{1/3+o(1)}$ bound [sweep]. Weighted Fermat searches go back to Lawrence
\cite{lawrence1895} (UNVERIFIED: cited through the sweep, not checked). The weighted sums
$Bp+Aq$ of the present paper are Lehman's mechanism; we claim nothing new about them.

\paragraph{McKee 1999.}
McKee \cite{mckee1999} speeds up Fermat's method using a prime $m$ dividing the square-root
variable, which restricts it modulo $m^2$; the resulting $N^{1/4+\eta}$ method is labelled
heuristic by the author, and its \S9 points back to McKee--Pinch \S5 [sweep].

\paragraph{McKee--Pinch 1996, \S5: a direct precursor.}
We read \S5 (pp.~221--223) of \cite{mckeepinch1996} from the public scan indexed at
\url{https://cr.yp.to/bib/1996/mckee-deterministic.html}. It contains the following.
\begin{itemize}
\item \emph{Lehmer's speed-up:} if $p\equiv q\equiv1\pmod k$, Fermat's method runs in
  $O(q/k^2)$.
\item \emph{Their extension:} for $p\equiv\pm q\pmod k$, Fermat runs in $O(q/k^{2-\epsilon})$,
  writing $n=u^2-v^2$ with $u$ or $v$ determined modulo $k^2$ from square roots modulo $k$.
  This is a modulus-squared lock.
\item \emph{Multipliers:} if $\lambda p\equiv\pm\mu q\pmod k$, the same method splits
  $\lambda\mu n$ in $O(\max(\lambda p,\mu q)/k^2\cdot n^{\epsilon})$.
\item \emph{A deterministic $O(n^{1/3+\epsilon})$ algorithm}, with $k=2\lfloor n^{1/3}\rfloor+1$
  and a pigeonhole lemma guaranteeing $\lambda\mu\le k$ with $\lambda p\equiv\pm\mu q\pmod k$.
\item \emph{A heuristic balanced method}, described by the authors as ``not however an
  algorithm'': take primes $k\in(n^{1/4},2n^{1/4})$ and small $\lambda,\mu$, and hope that
  $\lambda p-\mu q\equiv0\pmod k$ for some of them. Each $(\lambda\mu,k)$ costs $O(n^{\epsilon})$,
  and the expected time $O(n^{1/4+\epsilon})$ rests on the heuristic that $q-p$ has a prime
  factor in the band with probability $\gg1/\log n$.
\end{itemize}
(Here $n$ is their notation for the number to be factored.) So the core event of the present
paper (a band prime divides a small-weight combination $Bp-Aq$, with $N$-determined residues
and a modulus-squared reduction) is theirs. Their per-$k$ cost is $O(n^\epsilon)$ because
$q/k^2=O(1)$. We instead trade a larger window for fewer, smaller primes through one baby
list per band, and we prove the coverage statement. The contributions (i)--(v) of the
introduction are what we add.

\paragraph{Hittmeir 2018 and 2021.}
Hittmeir \cite{hittmeir2018} combines baby-step/giant-step with sieving of $p+q$ modulo many
small primes and proves $N^{1/4}\exp(-\Omega(\log N/\log\log N))$ up to polylogarithmic
factors [sweep]. Hittmeir \cite{hittmeir2021} gives a time--space trade-off for Lehman's
method, adapting Strassen-type polynomial collision detection to Lehman's weighted sums, and
proves $N^{2/9+o(1)}$ [sweep]. Modular-hyperbola sieving over many small primes is therefore
old; what differs here is that each prime is used \emph{selectively} (only the at most two
classes reached by a short weight vector) rather than to cover all classes.

\paragraph{St\u{a}nic\u{a}--Mulder--Hittmeir 2026.}
\cite{smh2026} studies unions of weighted modular-hyperbola sieve sets, their CRT behaviour,
and elliptic-trace obstructions [sweep]. This is closely related local arithmetic. As far as
the sweep found, it does not contain a global coverage theorem over prime pairs (our
Lemma~\ref{lem:sieve} summed over $q$) or a shared cofactor-powered search.

\paragraph{Harvey 2021; Harvey--Hittmeir 2022.}
Harvey \cite{harvey2021} proves deterministic complexity $O(N^{1/5}\log^{16/5}N)$ with a
single shared baby-step/giant-step search over Lehman's multiplier family, using
congruences $\alpha^{aq+bp}\equiv\alpha^{aN+b}\pmod p$ and multipoint evaluation over
$\Z/N\Z$ (Alg.~4.1--4.2) [sweep]. Harvey--Hittmeir \cite{hh2022} save a factor
$(\log\log N)^{3/5}$; their Proposition~3.3 contains the congruence $aq-bp\equiv0\pmod m$
and the modulus-squared lock that we use as Lemma~\ref{lem:lock}, with a single modulus
and lattice-chosen weights [sweep]. Harvey also remarks that exponent $1/6$ would need a
``fully square-root speedup'' of Lehman's search with $r\asymp N^{1/3}$ [sweep].

\paragraph{Gao--Feng--Hu--Pan 2025.}
For balanced semiprimes, \cite{gao2025} claims deterministic
$O(N^{1/5}\log^{13/5}N/(\log\log N)^{3/5})$ using rank-3 lattices [sweep; statements checked,
proofs not reproduced]. The exponent remains $1/5$.

\paragraph{Umans--Wang 2025; He--Sahai 2026.}
Umans and Wang \cite{umanswang2025} show that a Strong Prefactored Divisor Conjecture would
give deterministic exponent $1/6$ (their Thms.~5.2 and~5.5). He and Sahai
\cite{hesahai2026} refute the one-dimensional arithmetic-progression version at the proposed
parameters $(1/3,1/3)$; higher-rank versions remain open [sweep]. Our result is
unconditional but average-case. The covering statement we would need for a worst-case
version (Conjecture~\ref{conj:DCH} below) is a different object: we know of no implication
in either direction, and no transfer of the He--Sahai argument to it, but we have not
proved that none exists.

\paragraph{Oznovich--Volk 2025; Harvey--Hittmeir 2026.}
These give deterministic algorithms that find an element of large multiplicative order
modulo $N$ or a factor: Oznovich--Volk \cite{oznovichvolk2025} for target order
$D\ge N^{1/6}$ in time $D^{1/2+o(1)}$, and Harvey--Hittmeir \cite{hh2026} for every $D$
[sweep]. We do not use them. They produce an arbitrary element of large order, not one in
the prescribed image $2^{L}$, and large order modulo $N$ does not imply large order modulo
a specific prime $q$. Our clean-band argument avoids the issue.

\paragraph{Carella, arXiv:1308.2891.}
According to its abstract, Carella \cite{carella} presents a deterministic algorithm for
factoring a balanced integer $N=pq$ in $O(N^{1/6+\varepsilon})$ arithmetic operations, and
states equivalently that $N$ is factored in polynomial time given the least $(\log N)/6$ bits
of a factor; the abstract states no average-case restriction. We have not verified this
claim, and we take no position on it. The GPT-6.1-sol referee of our round 13 raised a
specific objection to the small-root argument in its \S5 (a lost primitive-polynomial
normalisation, after which the cited small-root condition would no longer support the
stated bounds). We have not checked that objection ourselves, so we do not adopt it. The
claim should be regarded as an unverified prior statement of the same exponent, by a
different (small-root) method.

\paragraph{The large sieve.}
We use the arithmetic large sieve in Montgomery's form \cite{montgomery1968} with the
Montgomery--Vaughan constant \cite{mv1973}; in fact only its prime terms are needed, and we
derive them from the additive inequality in Lemma~\ref{lem:sieve}.

\paragraph{Summary of the literature search.}
Two independent GPT-6.1-sol sweeps and our own searches found no established,
unconditional, oracle-free deterministic result below exponent $1/5$ that factors almost
all bounded-ratio prime pairs (Carella's claim above is unverified). Almost-all results that
rely on totient or elliptic-curve oracles exist but do not take only $N$ as input [sweep].
This is the outcome of a search, not a proof of priority.

%======================================================================
\section{The algorithm}\label{sec:alg}
%======================================================================

\subsection{Parameters}
All parameters depend on $n$ and $\varepsilon$ only, never on $N$. Constant-factor
conventions (floors and ceilings) do not affect the proofs.
\begin{itemize}
\item Weight height $H=\lceil n\,2^{\varepsilon n/2}\rceil$, and weight set
  \[V=\{(A,B)\in[1,H]^2:\gcd(A,B)=1,\ A\ne B\}.\]
  For large $H$, $H^2/2\le|V|\le H^2$.
\item Band scale $Y=\lceil 2^{n/6}H^{-1/3}\rceil$.
\item Four bands, $k=0,1,2,3$:
  \[\Pk=\{\ell\ \text{prime}:\ \ell\equiv3\pmod 4,\ 2^kY<\ell\le2^{k+1}Y\},\qquad
  L_k=\prod_{\ell\in\Pk}\ell^2.\]
\item Window bound $\spanbar=\lfloor\tfrac52H(2^{n/2}+1)\rfloor+2H+3$ and
  $\mbar=\lfloor\spanbar/(2^kY)^2\rfloor+2$.
\item Base $\beta=2$.
\end{itemize}
We record the scale separations used below. As $n\to\infty$ with $\varepsilon\le1/60$
fixed:
\begin{equation}\label{eq:scales}
  Y>H^2,\qquad 16Y<P/2,\qquad Y^4>\sqrt2\,P.
\end{equation}
The first needs $2^{n/6}>n^{7/3}2^{7\varepsilon n/6}$ (true for $\varepsilon<1/7$); the
third needs $2^{n/6}>\sqrt2\,H^{4/3}$ (true for $\varepsilon<1/4$).

\subsection{Steps}
Input: $N$ with $n$ bits.
\begin{enumerate}[label=\arabic*.]
\item If $N<4$ output \texttt{FAIL}. If $N$ is even output $2$. Put $a_0=\lceil\sqrt{N/2}\rceil$
  and $b_0=\lfloor\sqrt N\rfloor$; every $(p,q)\in\Sn$ has $p\in[a_0,b_0]$.
\item For $v=(A,B)\in V$ let $F_v(x)=Bx+AN/x$, $\mathrm{Lo}_v=\lfloor2\sqrt{ABN}\rfloor$ and
  $\mathrm{Hi}_v=\max_{x\in\{a_0,b_0\}}\lceil F_v(x)\rceil$. By AM--GM, $\mathrm{Lo}_v\le F_v(x)$
  for all $x>0$, and by convexity $F_v\le\mathrm{Hi}_v$ on $[a_0,b_0]$. Also
  $\mathrm{Hi}_v-\mathrm{Lo}_v\le(B+\sqrt2A)\sqrt N+2H+2\le\spanbar$.
\end{enumerate}
Then, for each block $k=0,1,2,3$:
\begin{enumerate}[label=\arabic*.,resume]
\item \emph{Exponent tree.} Build the balanced integer product tree of $\{\ell^2:\ell\in\Pk\}$.
  Write $\Pi(\text{node})$ for the product of the $\ell^2$ at the leaves below a node. Push $2$
  down the tree: a node holding $2^{L_k/\Pi(\text{node})}$ passes
  $2^{L_k/\Pi(\text{node})\cdot\Pi(\text{sibling})}$ to each child. This yields
  $u_\ell=2^{L_k/\ell^2}\bmod N$ for all $\ell\in\Pk$. Put $\gamma=u_{\ell_0}^{\ell_0^2}=2^{L_k}\bmod N$
  for any $\ell_0\in\Pk$. If some $\ell\in\Pk$ divides $N$ and $1<\ell<N$, output $\ell$.
\item \emph{Targets.} For $\ell\in\Pk$ and $v\in V$: let $a=ANB^{-1}\bmod\ell$. If
  $a^{(\ell-1)/2}\not\equiv1\pmod\ell$, skip. Otherwise $r_0=a^{(\ell+1)/4}\bmod\ell$, and the
  roots of $Br^2\equiv AN\pmod\ell$ are $r_0$ and $\ell-r_0$. For each root $r$ put
  \[
    s=Nr^{-1}\bmod\ell^2,\quad c=(Br+As)\bmod\ell^2,\quad
    g_{\mathrm{lo}}=\Bigl\lceil\frac{\mathrm{Lo}_v-c}{\ell^2}\Bigr\rceil,\quad
    g_{\mathrm{hi}}=\Bigl\lfloor\frac{\mathrm{Hi}_v-c}{\ell^2}\Bigr\rfloor,
  \]
  and $W_\tau=g_{\mathrm{hi}}-g_{\mathrm{lo}}+1\le\spanbar/\ell^2+1\le\mbar$. The triple
  $\tau=(v,\ell,r)$ is a \emph{target}. Let $T_k\le2|V||\Pk|$ be the number of targets.
\item \emph{Baby length.} Choose $m$ either as $m=\mbar$, or by the
  \emph{endpoint-competitive} rule: minimise $m+\sum_\tau\lceil W_\tau/m\rceil$ over an integer
  geometric grid whose last point is $W_{\max}=\max_\tau W_\tau\le\mbar$ (if $T_k=0$ the block
  is skipped). Under either rule
  $m\le\mbar$ and $m+T'\le\mbar+T_k$, where $T'=\sum_\tau\lceil W_\tau/m\rceil$.
\item \emph{Giants.} Precompute $w_\ell=u_\ell^N$. For each target and $0\le j<\lceil W_\tau/m\rceil$:
  \[
    G_{\tau,j}=w_\ell^{A}\,u_\ell^{\,B-c+\ell^2}\,\gamma^{-(g_{\mathrm{lo}}+1)}\,\gamma^{-jm}
    \equiv u_\ell^{\,B+AN-c}\,\gamma^{-(g_{\mathrm{lo}}+jm)}\pmod N,
  \]
  using $u_\ell^{\ell^2}=\gamma$; $\gamma$ is invertible because $N$ is odd.
\item \emph{Collision detection.} Build $f(x)=\prod_{0\le i<m}(x-\gamma^i)\in(\Z/N\Z)[x]$ by a
  product tree. Evaluate $f$ at all giants by a remainder tree; only monic divisors occur,
  so this is valid over $\Z/N\Z$. For each giant compute $\gcd(f(G),N)$; if it lies strictly
  between $1$ and $N$, output it.
\end{enumerate}
If no block outputs a divisor, output \texttt{FAIL}.

\begin{remark}\label{rem:alg}
(a) No multiplicative order is computed, assumed or certified, and no randomness is used.
(b) The symmetric weight $A=B$ is excluded: by Lemma~\ref{lem:modq} its true giant always
collides modulo $q$ too, so it can never give a one-sided collision.
(c) The four bands are used only in Lemma~\ref{lem:clean}.
(d) The condition $m\le\mbar$ suffices for the success proof, but not for the cost bound:
$m=1$ on adversarial inputs would make $T'$ far too large. That is why one of the two rules
in step~5 is mandatory.
(e) An implementation (Python, gmpy2 and python-flint) is in
\texttt{src/dcltheorem/algorithm.py}. Every candidate output $x$ passes a check $1<x<N$,
$x\mid N$; an exhaustive regression test covers all $2\le N\le10^4$ for two parameter
designs.
\end{remark}

%======================================================================
\section{Proof of \texorpdfstring{Theorem~\ref{thm:main}}{Theorem 1.1}}\label{sec:proofs}
%======================================================================

Throughout this section fix $(p,q)\in\Sn$, $N=pq$, a block $k$, and write $L=L_k$,
$u=u_\ell$, $d=\ord_q(2)$. Since $p>P/2>16Y$ by \eqref{eq:scales}, every band prime is smaller
than both factors.

\subsection{Modulus lock, coverage and windows}

\begin{lemma}[Modulus-squared lock; {\cite[Prop.~3.3]{hh2022}}]\label{lem:lock}
Let $\ell\in\Pk$, $(A,B)\in V$, and let $r$ be a root of $Br^2\equiv AN\pmod\ell$. If
$p\equiv r\pmod\ell$, then $g:=(Bp+Aq-c)/\ell^2\in\Z$.
\end{lemma}
\begin{proof}
Since $\ell>Y>H\ge A,B$ and $\ell<p<q$, we have $\ell\nmid ABN$. Write $p=r+\ell x$. Modulo
$\ell^2$, $p^{-1}\equiv r^{-1}(1-\ell xr^{-1})$, so $q=Np^{-1}\equiv s-\ell x\,sr^{-1}$. Hence
$Bp+Aq\equiv c+\ell x(B-Asr^{-1})\pmod{\ell^2}$. Now $sr^{-1}\equiv Nr^{-2}\pmod\ell$ and
$Br^2\equiv AN$, so $B-ANr^{-2}\equiv0\pmod\ell$, and the last term vanishes modulo
$\ell^2$.
\end{proof}

\begin{lemma}[Coverage criterion]\label{lem:coverage}
$p\bmod\ell$ is one of the computed roots if and only if $\ell\mid Bp-Aq$. In that case the
Euler test in step~4 passes and the target $\tau=(v,\ell,p\bmod\ell)$ exists; we say that
$\tau$ \emph{covers} $p$.
\end{lemma}
\begin{proof}
$Bp^2-AN=p(Bp-Aq)$ and $\ell\nmid p$, so $Bp^2\equiv AN\pmod\ell$ iff $\ell\mid Bp-Aq$. If so,
$a\equiv p^2\not\equiv0$ is a quadratic residue, and for $\ell\equiv3\pmod4$ the two roots
are $\pm a^{(\ell+1)/4}$.
\end{proof}

\begin{lemma}[Window]\label{lem:window}
If $\tau$ covers $p$, then $g_{\mathrm{lo}}\le g\le g_{\mathrm{hi}}$. Put $i_{\mathrm{full}}=g-g_{\mathrm{lo}}$,
$j^*=\lfloor i_{\mathrm{full}}/m\rfloor$ and $i^*=i_{\mathrm{full}}-j^*m\in[0,m)$. Then
$j^*<\lceil W_\tau/m\rceil$, so $G_{\tau,j^*}$ is one of the computed giants; we call it the
\emph{true giant}.
\end{lemma}
\begin{proof}
$p\in[a_0,b_0]$, so $\mathrm{Lo}_v\le F_v(p)=Bp+Aq\le\mathrm{Hi}_v$. Divide by $\ell^2$ after
subtracting $c$. Finally $i_{\mathrm{full}}\le W_\tau-1$.
\end{proof}

\subsection{The collision modulo \texorpdfstring{$p$}{p} is exact}

\begin{lemma}[Mod $p$]\label{lem:modp}
If $\tau$ covers $p$, then $G_{\tau,j^*}\equiv\gamma^{i^*}\pmod p$. No hypothesis on any
multiplicative order is needed.
\end{lemma}
\begin{proof}
We have the identity
\[
  B+AN-(Bp+Aq)=(p-1)(Aq-B),
\]
so $B+AN-c=\ell^2g+(p-1)(Aq-B)$. Since $2^{p-1}\equiv1\pmod p$,
$u^{B+AN-c}\equiv2^{(L/\ell^2)\ell^2g}=\gamma^g\pmod p$. Multiply by
$\gamma^{-(g_{\mathrm{lo}}+j^*m)}$.
\end{proof}

\begin{lemma}[Mod $q$]\label{lem:modq}
If $\tau$ covers $p$, then $G_{\tau,j^*}\equiv\gamma^{i^*}u^{(A-B)(p-1)}\pmod q$.
\end{lemma}
\begin{proof}
Modulo $q-1$ we have $N\equiv p$, so $B+AN-c\equiv B+Ap-c$; also
$\ell^2g=Bp+Aq-c\equiv Bp+A-c$. The difference is $(A-B)(p-1)$. Multiplying by $L/\ell^2$
preserves the congruence modulo $q-1$, and $2^{q-1}\equiv1\pmod q$.
\end{proof}

\begin{definition}
The true giant $G=G_{\tau,j^*}$ is \emph{FULL} if $G\equiv\gamma^{i'}\pmod q$ for some
$i'\in[0,m)$.
\end{definition}

\begin{lemma}[One-sided extraction]\label{lem:extract}
If $\tau$ covers $p$ and its true giant $G$ is not FULL, then $\gcd(f(G),N)=p$.
\end{lemma}
\begin{proof}
$f(G)=\prod_{i<m}(G-\gamma^i)$. The factor $i=i^*$ vanishes modulo $p$ by
Lemma~\ref{lem:modp}. Modulo $q$ no factor vanishes, by definition, and $\Z/q\Z$ is a field.
\end{proof}

\subsection{The clean band and the FULL criterion}

\begin{lemma}[Clean band]\label{lem:clean}
Call block $k$ \emph{clean for $q$} if no $\ell\in\Pk$ divides $q-1$. At least one of the four
blocks is clean for $q$.
\end{lemma}
\begin{proof}
Every band prime exceeds $Y$. If all four blocks were unclean, $q-1$ would have four
distinct prime factors exceeding $Y$, so $q-1>Y^4>\sqrt2P>q$ by \eqref{eq:scales}, a
contradiction.
\end{proof}

\begin{lemma}[FULL on a clean block]\label{lem:full}
If block $k$ is clean for $q$, then the true giant is FULL if and only if
\[
  d\mid(A-B)(p-1)-\ell^2\delta\quad\text{for some }\delta\in[-i^*,\,m-1-i^*]\subset(-m,m).
\]
\end{lemma}
\begin{proof}
By Lemma~\ref{lem:modq}, FULL holds iff $u^{(A-B)(p-1)}\equiv\gamma^{\delta}\pmod q$ for some
$\delta=i'-i^*$, i.e.\ iff $2^{(L/\ell^2)[(A-B)(p-1)-\ell^2\delta]}\equiv1\pmod q$. Since $d\mid q-1$
and $\gcd(L,q-1)=1$ on a clean block, $\gcd(d,L/\ell^2)=1$, and the factor $L/\ell^2$ can be
cancelled.
\end{proof}

The exact asymmetric range of $\delta$ matters for the equivalence; below we only use the
enlarged range $|\delta|<\mbar$ as a superset for counting.

\subsection{Three counting lemmas}

\begin{lemma}[Small orders are rare]\label{lem:smallorder}
For $Q_0\ge2$ and $T\ge1$,
\[
  \#\{q>Q_0\ \text{prime}:\ \ord_q(2)\le T\}\le\frac{T(T+1)}{2\log_2Q_0}.
\]
\end{lemma}
\begin{proof}
Such a $q$ divides $2^j-1$ for some $j\le T$. Since $2^j-1<2^j$, it has fewer than
$j/\log_2Q_0$ prime factors exceeding $Q_0$. Sum over $j\le T$.
\end{proof}

\begin{lemma}[Coverage sieve]\label{lem:sieve}
Fix $q$ and $k$. Let $U_k(q)$ be the set of integers $p\in[P/2,P]$ such that
$\ell\nmid Bp-Aq$ for all $(v,\ell)\in V\times\Pk$. For $n\ge n_0$,
\[
  |U_k(q)|\le\frac{P/2+1+4^{k+1}Y^2}{J_k},\qquad
  J_k:=\sum_{\ell\in\Pk}\frac{|V|}{\ell-|V|}\ \ge\ |V|\sum_{\ell\in\Pk}\frac1\ell\ \ge\ \frac{c_1H^2}{n},
\]
with an absolute constant $c_1>0$.
\end{lemma}
\begin{proof}
For $\ell\in\Pk$ we have $\ell\nmid q$, so $\ell\mid Bp-Aq$ iff $p\equiv AqB^{-1}\pmod\ell$. These
$|V|$ classes are distinct: $AqB^{-1}\equiv A'qB'^{-1}$ would give $\ell\mid AB'-A'B$, but
$0<|AB'-A'B|<H^2<Y<\ell$. So $U=U_k(q)$ avoids $\omega:=|V|<\ell$ residue classes modulo
every $\ell\in\Pk$.

Write $e(x)=\exp(2\pi ix)$. Let $S(\alpha)=\sum_{u\in U}e(u\alpha)$ and $Z(h)=\#\{u\in U:u\equiv h\pmod\ell\}$. Then
$\sum_{a=0}^{\ell-1}|S(a/\ell)|^2=\ell\sum_hZ(h)^2\ge\ell|U|^2/(\ell-\omega)$ by Cauchy--Schwarz,
since $Z$ is supported on $\ell-\omega$ classes. Removing $a=0$,
\[
  \sum_{a=1}^{\ell-1}|S(a/\ell)|^2\ge\frac{\omega}{\ell-\omega}|U|^2.
\]
The points $a/\ell$ with $\ell\in\Pk$, $1\le a<\ell$, are $Q^{-2}$-separated modulo $1$ with
$Q=2^{k+1}Y$, and $U$ lies in an interval of at most $N_0=P/2+1$ integers. The additive
large sieve inequality \cite{mv1973} gives
$\sum_{\ell,a}|S(a/\ell)|^2\le(N_0+Q^2)|U|$. Combining,
$|U|\,J_k\le N_0+Q^2$, since $\omega=|V|$. Finally $\omega/(\ell-\omega)\ge|V|/\ell$.

By the prime number theorem in the progression $3\bmod4$ (effective, the modulus being
fixed), $\sum_{\ell\in\Pk}1/\ell\ge c'/n$ for an absolute $c'>0$, and $|V|\ge H^2/2$ since
coprime pairs have density $6/\pi^2$. So $J_k\ge c'H^2/(2n)$, i.e.\ $c_1=c'/2$.
\end{proof}

\begin{lemma}[FULL is rare]\label{lem:crt}
Fix $q$ with $d=\ord_q(2)>T$ and a block $k$ clean for $q$. Let $F_k(q)$ be the set of
integers $p\in[P/2,P]$ for which some $(v,\ell)\in V\times\Pk$ has $\ell\mid Bp-Aq$ and
$d\mid(A-B)(p-1)-\ell^2\delta$ for some $|\delta|<\mbar$. Then
\[
  |F_k(q)|\le2|V||\Pk|\,\mbar\Bigl(\frac{PH}{2^{k+1}YT}+1\Bigr).
\]
\end{lemma}
\begin{proof}
Fix $(v,\ell,\delta)$ and let $g'=\gcd(A-B,d)\le H$ (recall $A\ne B$ and $|A-B|<H$). The
congruence $(A-B)(p-1)\equiv\ell^2\delta\pmod d$ is either unsolvable or holds on exactly
one class of $p$ modulo $d/g'$. Cleanliness gives $\ell\nmid q-1$, hence $\ell\nmid d$, so
together with $p\equiv AqB^{-1}\pmod\ell$ the CRT gives one class modulo $\ell d/g'\ge\ell T/H$.
An interval of length $P/2$ contains at most $PH/(2\ell T)+1$ of its elements, and
$\ell>2^kY$. There are at most $|V||\Pk|(2\mbar-1)$ triples.
\end{proof}

\subsection{The exceptional set}
Let $T=\lfloor2^{(1/4-\varepsilon)n}\rfloor$. For $(p,q)\in\Sn$ let $k(q)$ be the least block
that is clean for $q$ (Lemma~\ref{lem:clean}).

\begin{proposition}[Success certificate]\label{prop:cert}
If (a) $d=\ord_q(2)>T$, (b) $p\notin U_{k(q)}(q)$, and (c) $p\notin F_{k(q)}(q)$, then
$\DCL$ outputs a proper divisor of $N$.
\end{proposition}
\begin{proof}
By (b) and Lemma~\ref{lem:coverage}, some target $\tau$ of block $k(q)$ covers $p$; by
Lemma~\ref{lem:window} its true giant is computed. By (c) and Lemma~\ref{lem:full} it is
not FULL, because the exact $\delta$-range lies inside $|\delta|<m\le\mbar$. By
Lemma~\ref{lem:extract} its gcd is $p$, so block $k(q)$ outputs a proper divisor unless an
earlier block already did. Either output is a proper divisor; it may be $q$ rather than
$p$.
\end{proof}

Hence the \texttt{FAIL} pairs lie in $E_a\cup E_b\cup E_c$, where
$E_a=\{d\le T\}$, $E_b=\{d>T,\ p\in U_{k(q)}(q)\}$, $E_c=\{d>T,\ p\in F_{k(q)}(q)\}$. We
sum over the at most $3P/n$ primes $q\in(P/\sqrt2,\sqrt2P)$ and divide by $|\Sn|\ge c_0P^2/n^2$.
This normalisation (one fixed interval for $p$, integer survivors counted, then summed over
prime $q$) needs no prime counts in short intervals.

\paragraph{$E_a$.} With $Q_0=2^{(n-1)/2}<q$, Lemma~\ref{lem:smallorder} bounds the bad $q$
by $T(T+1)/(n-1)$, each with at most $P/2+1$ partners. So
$|E_a|\le(1+o(1))2^{(1-2\varepsilon)n}/n$ and $|E_a|/|\Sn|\le(n/c_0)2^{-2\varepsilon n}(1+o(1))$.

\paragraph{$E_b$.} By Lemma~\ref{lem:sieve} and $Y^2=o(P)$,
$\sum_q|U_{k(q)}(q)|\le(3P/n)\cdot(P/2+1+256Y^2)\cdot n/(c_1H^2)\le CP^2/H^2$. So
$|E_b|/|\Sn|\le Cn^2/H^2\le C2^{-\varepsilon n}$.

\paragraph{$E_c$.} Use $|V|\le H^2$, $|\Pk|\le2^kY$ and $\mbar\le3HP/(4^kY^2)$ in
Lemma~\ref{lem:crt}: $|F_k(q)|\le6H^4P^2/(Y^2T)+6H^3P/Y$. Summing over $q$ and using
$Y^2\ge P^{2/3}H^{-2/3}$,
\[
  \frac{|E_c|}{|\Sn|}\le\frac{3P}{n}\Bigl(\frac{6H^4P^2}{Y^2T}+\frac{6H^3P}{Y}\Bigr)\frac{n^2}{c_0P^2}
  \le Cn\Bigl(\frac{H^{14/3}P^{1/3}}{T}+\frac{H^{10/3}}{P^{1/3}}\Bigr).
\]
With $H\le2n2^{\varepsilon n/2}$, $P=2^{n/2}$ and $T\asymp2^{(1/4-\varepsilon)n}$, the two terms are
at most $n^{O(1)}2^{(-1/12+10\varepsilon/3)n}$ and $n^{O(1)}2^{(-1/6+5\varepsilon/3)n}$. For every
$\varepsilon\le1/60$ we have $-1/12+10\varepsilon/3\le-\varepsilon-1/90$ (equivalently $13\varepsilon/3\le13/180$), so
the first term is at most $n^{O(1)}2^{-n/90}\cdot2^{-\varepsilon n}=o(2^{-\varepsilon n})$; the second is
smaller still. (Any fixed $\varepsilon<1/52$ would do here, since that is what
$-1/12+13\varepsilon/3<0$ requires.)

\paragraph{Conclusion.} The total exceptional fraction is $O(2^{-\varepsilon n})=O(N^{-\varepsilon})$,
which proves part (iii). Part (i) holds because only gcds in $(1,N)$ and band primes
$1<\ell<N$ dividing $N$ are output, and inputs $N<4$ return \texttt{FAIL}. \qed

\begin{remark}[Where the constants come from]
The bound $\varepsilon\le1/60$ is used only in $E_c$ (which needs $\varepsilon<1/52$) and in the scale
separations \eqref{eq:scales} ($\varepsilon<1/7$ and $\varepsilon<1/4$). The binding term is $E_b$: the
large sieve gives $n^2/H^2$, so an exceptional fraction $N^{-\varepsilon}$ costs $H=N^{\varepsilon/2}$,
i.e.\ a factor $N^{5\varepsilon/6}$ in the running time. A Poisson heuristic predicts an
uncovered fraction $e^{-\lambda}$ with $\lambda\asymp H^2/n$, which would give $N^{-\varepsilon}$ already
at polylogarithmic $H$ (conjectured; Section~\ref{sec:open}).
\end{remark}

\subsection{Bit complexity}\label{sec:cost}
We assume quasi-linear integer multiplication $M(b)=O(b\log b)$ \cite{hvdh2021}
(Sch\"onhage--Strassen costs an extra $\log\log$). Per block, with $T_k$ targets and
$m\le\mbar$:

\begin{center}
\small
\begin{tabular}{p{0.42\textwidth}p{0.22\textwidth}p{0.24\textwidth}}
\toprule
step & operations & bit cost\\
\midrule
prime sieve to $16Y$; $V$; $\mathrm{Lo}_v,\mathrm{Hi}_v$ & $O(Y\log\log Y)$; $O(H^2)$ square roots & $\tilde O(Y+H^2n)$\\
integer product tree of $\{\ell^2\}$ ($\log_2L_k\le3\cdot2^kY$) & balanced tree & $O(M(Y)\log Y)$\\
exponent tree; $\gamma,\gamma^{-1}$ & $O(Y\log Y)$ mults mod $N$, 1 inversion & $O(Y\log Y\,M(n))$\\
roots, $s$, $c$, windows & $O(|V||\Pk|\log Y)$ mults mod $\ell,\ell^2$ & $\tilde O(H^2Y)$\\
$w_\ell=u_\ell^N$; giants & $O(n|\Pk|+nT_k+T')$ mults mod $N$ & $O((n(T_k+Y)+T')M(n))$\\
choice of $m$; small inversions and reductions & $O(nT_k)$ small-integer ops & $\tilde O(T_kn^{O(1)})$\\
babies $\gamma^i$ & $m$ mults & $O(mM(n))$\\
$f$ by product tree; remainder tree at $T'$ giants & $O((M(m)+M(T'))\log(m+T'))$ ring ops & $\tilde O((m+T')n)$\\
gcds & $T'$ gcds of $n$-bit integers & $O(T'M(n)\log n)$\\
\bottomrule
\end{tabular}
\end{center}

Polynomial products over $\Z/N\Z$ use Kronecker substitution with slots wider than
$2n+\lceil\log_2(D+1)\rceil$, so a degree-$D$ product costs $O(M(D(2n+\log D)))$ bit
operations \cite{vzgg}. Remainders by monic polynomials use reversal and Newton inversion of
a power series with constant term $1$, so no inverse of a nonunit is ever needed.

By step~5, $m+T'\le\mbar+T_k\le\tfrac52HP/Y^2+4+2H^2\cdot2^{k+1}Y$. Over the four blocks the
total is
\[
  O\bigl(n^{O(1)}(HP/Y^2+H^2Y)\bigr)=O\bigl(n^{O(1)}H^{5/3}P^{1/3}\bigr)
  =O\bigl(n^{O(1)}2^{(1/6+5\varepsilon/6)n}\bigr),
\]
because $Y\asymp P^{1/3}H^{-1/3}$. This proves part (ii). \qed

\subsection{Proof of \texorpdfstring{Corollary~\ref{cor:polylog}}{Corollary 1.2}}
Take $H=\lceil n^{1+a}\rceil$, $Y$ as before, and in the analysis $T=\lfloor2^{n/4}n^{-1-a}\rfloor$. The
scale separations \eqref{eq:scales} hold for large $n$. The three exceptional terms become
$O(n^{-1-2a})$ ($E_a$), $O(n^{-2a})$ ($E_b$) and
$n^{O(1)}\cdot2^{-n/12}\cdot n^{O(1+a)}=o(n^{-2a})$ ($E_c$). The cost is
$n^{O(1)}H^{5/3}2^{n/6}=N^{1/6}(\log N)^{O(1+a)}$. \qed

%======================================================================
\section{Context: a \texorpdfstring{$1/6$}{1/6} floor for full-candidate covering}\label{sec:lower}
%======================================================================

This section is context, not part of the proof of Theorem~\ref{thm:main}. It explains why
many moduli are needed to go below $1/5$ in Lehman-type searches, and why $1/6$ is the floor
for one specific notion of covering. \textbf{Quantifier warning.} Every result here concerns
families that cover \emph{every integer candidate} in the balanced range for a \emph{fixed}
$N$. None of them bounds an algorithm that, like ours, only needs to reach the true factor for
most inputs, and none is a lower bound for factoring.

The results were obtained in the project notes (\texttt{notes/cycloharvey\_results.md} \S3,
\texttt{notes/multimod\_results.md} \S3) and checked by one independent AI critic and by the
GPT-6.1-sol round-12 referee, who judged Theorems~\ref{thm:A}, \ref{thm:Aprime}
and~\ref{thm:C} correct in the stipulated model and Theorem~\ref{thm:B} correct given a
short counting argument. Below, Lemma~\ref{lem:K}, Theorem~\ref{thm:C} and the odd-prime case
of Lemma~\ref{lem:classcount} are proved in full; Lemma~\ref{lem:C} and the prime-$2$ case of
Lemma~\ref{lem:classcount} are given as sketches with references to the notes.

\subsection{The model}
Fix a degree bound $d_0\ge1$ and $\theta>1$. Write $R=\sqrt N$; the \emph{candidates} are the
integers $p\in[R/\sqrt\theta,R]$, and we put $t=p/R\in J=[\theta^{-1/2},1]$. A \emph{target} $\tau$
consists of:
\begin{itemize}
\item a modulus $M_\tau\ge1$ with $\gcd(M_\tau,N)=1$ and a \emph{unit} residue class
  $r_\tau\in(\Z/M_\tau\Z)^\times$;
\item a support $I_\tau\subseteq[-d_0,d_0]\setminus\{0\}$ with $K_\tau=|I_\tau|$, nonzero
  coefficients $\alpha_{\tau,i}\in\Z[1/M_\tau]$ ($i\in I_\tau$) and a constant
  $c_\tau\in\Z[1/M_\tau]$, defining the exponent map
  \[
    h_\tau(x)=\sum_{i\in I_\tau,\,i>0}\alpha_{\tau,i}x^{i}+\sum_{i\in I_\tau,\,i<0}\alpha_{\tau,i}(N/x)^{|i|}-c_\tau ,
  \]
  i.e.\ a polynomial in $p$ and $q=N/p$, shifted by a constant;
\item an integer window $[U_\tau,U_\tau+W_\tau)$ with $W_\tau$ a positive integer.
\end{itemize}
The \emph{integrality condition} is: for each prime $\ell\mid M_\tau$, $h_\tau(x)\in\Z_\ell$ for every
$x\in\Z_\ell$ with $x\equiv r_\tau\pmod{\ell^{v_\ell(M_\tau)}}$. Such $x$ are $\ell$-adic units (this is
where the unit-class hypothesis is used) and $N$ is an $\ell$-adic unit, so $h_\tau(x)$ is
defined. The condition is imposed on the whole residue class, as an identity; requiring it
only at the actual divisors of $N$ would make the model vacuous. It guarantees that if a
candidate $p\equiv r_\tau\pmod{M_\tau}$ divides $N$, then $h_\tau(p)\in\Z$.

The target $\tau$ \emph{covers} the candidate $p$ if $p\equiv r_\tau\pmod{M_\tau}$ and the real
number $h_\tau(p)$ (with $q$ read as $N/p$) lies in $[U_\tau,U_\tau+W_\tau)$. A family
\emph{covers} if every candidate is covered by some target. Its \emph{cost} is any $X$ with
\begin{enumerate}[label=(C\arabic*)]
\item at most $X$ distinct triples $(M_\tau,r_\tau,(\alpha_{\tau,i})_i)$;
\item at most $X$ targets, i.e.\ nonempty original windows;
\item $\sum_\tau W_\tau\le X^2$.
\end{enumerate}
These hold for a sumset baby-step/giant-step search with baby list length $m\le X$ and at
most $X$ giants, whether the list is shared, per modulus or mixed: every window costs
$\lceil W_\tau/m\rceil\ge1$ giants, so (C2) holds, and
$\sum_\tau W_\tau\le\sum_\tau m\lceil W_\tau/m\rceil\le X\max m\le X^2$. Call this model $\mathfrak M^+$.
With a single modulus it is the identity model in which Lehman, Hittmeir, Harvey and
Harvey--Hittmeir argue. The targets of $\DCL$ fit it with $M_\tau=\ell$, $r_\tau$ a root of
$Br^2\equiv AN$ (a unit, since $\ell\nmid AN$), $h_\tau=(Bp+Aq-c)/\ell^2$ and $W_\tau\approx HR/\ell^2$;
integrality is Lemma~\ref{lem:lock}.

\subsection{Per-target coverage}
With $p=Rt$ and $q=N/p=R/t$ we have $h_\tau(Rt)=G_\tau(t)-c_\tau$ with
$G_\tau(t)=\sum_{i\in I}b_it^i$ and $b_i=\alpha_{\tau,i}R^{|i|}$.

\begin{lemma}[Coverage of one map; sketch]\label{lem:C}
For any support $I$ with $0\notin I$, $|I|=K$, and any measurable $E\subset\R$,
$|\{t\in J:G(t)\in E\}|\le C_I(|E|/\|b\|_\infty)^{1/K}$.
\end{lemma}
\noindent\emph{Sketch} (full proof: \texttt{notes/cycloharvey\_results.md} \S3.2, Lemmas W
and C). The generalised Wronskian $\det[(i)_k]_{k\le K,\,i\in I}=\prod_{i\in I}i\prod_{i<i'}(i'-i)$ is
non-zero, so on $J$ the vector $(G',\dots,G^{(K)})$ is an invertible matrix times $b$, and
$J$ splits into $O_I(1)$ pieces on each of which some $|G^{(k)}|$, $k\le K$, is
$\ge c_I\|b\|_\infty$ and $G,G'$ are monotone. On each piece apply the van der Corput sublevel-set
estimate, after moving $E$ to an interval at the flat end; when $|E|>\|b\|_\infty$ the bound is
trivial.

\begin{lemma}[Denominators]\label{lem:K}
If $h_\tau$ satisfies the integrality condition, then $C_{d_0}M_\tau^{K_\tau}\alpha_{\tau,i}\in\Z$ for all
$i$, where $C_{d_0}$ depends only on $d_0$ and has only prime factors at most $2d_0$.
\end{lemma}
\begin{proof}
Coefficients are $\ell$-integral for $\ell\nmid M=M_\tau$ by definition. Fix $\ell^a\Vert M$, write
$r=r_\tau$, $K=K_\tau$, and let $v^*=\min_iv_\ell(\alpha_i)$. Since $r$ is an $\ell$-adic unit,
$f(y)=h_\tau(r+\ell^ay)$ is given on $\Z_\ell$ by a convergent series
$f(y)=T_0'+\sum_{n\ge1}T_n\ell^{an}y^n$ with
$T_n=\sum_i\alpha_iN^{\max(0,-i)}\binom{i}{n}r^{i-n}$ (generalised binomial coefficients for
$i<0$). Each $T_n$ is a $\Z_\ell$-combination of the $\alpha_i$, so $v_\ell(T_n)\ge v^*$. Moreover
$(T_1,\dots,T_K)=W(r)\alpha$ where, after removing unit factors, $\det W(r)$ is the Wronskian above
divided by $\prod_{n\le K}n!$; let $c_W=v_\ell(\det W)$, which is bounded in terms of $d_0$ and is
$0$ for $\ell>2d_0$. Hence $\min_{j\le K}v_\ell(T_j)\le v^*+c_W$; let $j_0\le K$ attain the minimum.

The $j_0$-th finite difference $\Delta^{j_0}f(0)=\sum_{n\ge j_0}T_n\ell^{an}j_0!\,S(n,j_0)$ (Stirling
numbers of the second kind) lies in $\Z_\ell$ because $f$ is $\Z_\ell$-valued. If $a>c_W$, the term
$n=j_0$ has valuation $v_\ell(T_{j_0})+aj_0+v_\ell(j_0!)\le v^*+c_W+aj_0+v_\ell(j_0!)$, while every term
with $n>j_0$ has valuation at least $v^*+a(j_0+1)+v_\ell(j_0!)$, which is larger. So
$v_\ell(T_{j_0})+aj_0+v_\ell(j_0!)\ge0$, and
$v^*\ge v_\ell(T_{j_0})-c_W\ge-aK-v_\ell(K!)-c_W$. If $a\le c_W$, apply the same argument to a subclass
modulo $\ell^{c_W+1}$ (integrality on the class implies it on subclasses); the loss is bounded
in terms of $d_0$. Since $v_\ell(K!)=0$ and $c_W=0$ for $\ell>2d_0$, the claim follows.
\end{proof}

\begin{proposition}[Per-target bounds]\label{prop:pertarget}
Let $\nu=\max_{i\in I_\tau}|i|$ be the degree of $\tau$. Merge all targets with the same triple
$(M,r,\alpha)$ and let $W$ be the total width of their windows and $w$ their number. The number
of candidates covered by the merged targets is at most:
\begin{enumerate}[label=(\roman*)]
\item if $\nu=1$, $K=2$: $C(RW/\eta)^{1/2}+Cw$, where we write
  $\alpha_1=B/(C_{d_0}M^2)$, $\alpha_{-1}=A/(C_{d_0}M^2)$ with integers $A,B$ (Lemma~\ref{lem:K}) and
  $\eta=\max(|A|,|B|)\ge1$;
\item if $\nu=1$, $K=1$: $CW+Cw$;
\item if $\nu\ge2$: $CR^{1-\nu/K}W^{1/K}+Cw$.
\end{enumerate}
In all cases the modulus $M$ cancels.
\end{proposition}
\begin{proof}
The candidates covered lie in $\{p\equiv r\pmod M\}$ and in the preimage of the union $E$ of the
(shifted) windows, which has measure at most $W$. The preimage of each window is a union of
at most $C_I$ intervals in $t$, so the number of integers $p\equiv r\pmod M$ with $p/R$ in the
preimage is at most $R\mu/M+C_Iw$, where $\mu$ is the measure of the preimage in $t$. By
Lemma~\ref{lem:C}, $\mu\le C(W/\|b\|_\infty)^{1/K}$. (i) $\|b\|_\infty=R\eta/(C_{d_0}M^2)$, so
$R\mu/M\le C(RW/\eta)^{1/2}$. (ii) $\alpha$ is a nonzero element of $(C_{d_0}M)^{-1}\Z$, so
$\|b\|_\infty\ge R/(C_{d_0}M)$ and $R\mu/M\le CW$. (iii) Some $|i|=\nu$ has $\alpha_i\ne0$, so
$\|b\|_\infty\ge R^\nu/(C_{d_0}M^{K})$ and $R\mu/M\le CR^{1-\nu/K}W^{1/K}$.
\end{proof}

Equality in (i) with $\eta=O(1)$ and $W\approx R/M^2$ means that a target covers its whole
residue class across the range, $R/M\approx(RW)^{1/2}$ candidates. That is what the targets of
$\DCL$ do.

\subsection{Statements}
Constants below depend on $d_0$ and $\theta$ only. In each proof, (C2) bounds the total rounding
term $C\sum w$ by $CX$.

\begin{theorem}[Theorem A; proved in the model]\label{thm:A}
Every covering family in $\mathfrak M^+$ for $N$ has cost $X\ge c_{d_0,\theta}N^{1/6}$.
\end{theorem}
\begin{proof}
There are $\asymp R$ candidates. Suppose $X\le\delta R^{2/5}$ with $\delta$ small. By
Proposition~\ref{prop:pertarget}(iii) and H\"older over at most $X$ merged triples
($\sum W^{1/K}\le X^{1-1/K}(\sum W)^{1/K}\le X^{1+1/K}$), the targets of degree $\nu\ge2$ cover at
most $C\sum_{(\nu,K)}R^{1-\nu/K}X^{1+1/K}+CX$ candidates; since $K\le2\nu$ and $\nu\ge2$ give
$5\nu\ge2K+2$, each term is at most $C\delta^{1+1/K}R$. By (ii), the degree-one single-term
targets cover at most $CX^2+CX\le C\delta^2R^{4/5}+CX=o(R)$. So for small $\delta$ the degree-one
two-term targets must cover a positive proportion of the candidates. By (i) and
Cauchy--Schwarz they cover at most
$CR^{1/2}\sum W^{1/2}+CX\le CR^{1/2}X^{1/2}(\sum W)^{1/2}+CX\le CR^{1/2}X^{3/2}+CX$ (using
$\eta\ge1$), which forces $X\ge cR^{1/3}=cN^{1/6}$. Otherwise $X>\delta R^{2/5}=\delta N^{1/5}\ge N^{1/6}$
for large $N$.
\end{proof}

\begin{theorem}[Theorem A$'$; proved in the model]\label{thm:Aprime}
If the degree-one targets of a covering family cover fewer than half of the candidates,
then $X\ge cN^{1/5}$, whatever the moduli. In particular, families whose targets all have
degree $\ge2$ never beat $N^{1/5}$.
\end{theorem}
\begin{proof}
Then the targets of degree $\ge2$ cover a positive proportion, and the first step of the
previous proof shows that this needs $X\ge cR^{2/5}$.
\end{proof}

\begin{lemma}[Class count]\label{lem:classcount}
Let $M\ge1$ with $\gcd(M,N)=1$ and let $A,B$ be nonzero integers. Let $\rho$ be the number of
unit classes $r\bmod M$ for which $(Bp+Aq)/M^2$ satisfies the integrality condition (for some
constant). Then $\rho\le4^{\omega(M)}\gcd(A,B,M)$. For $\nu=1$, $K=2$ the proof of
Lemma~\ref{lem:K} gives $c_W=0$ and $v_\ell(K!)=0$ at every odd $\ell$, so writing the
denominator as $M^2$ loses nothing there, and the proof below is complete for odd $M$. At
$\ell=2$ Lemma~\ref{lem:K} allows an extra bounded power of $2$ in the denominator; with it the
local count changes by a bounded factor. That last step is a sketch, checked by brute force
for composite $M\le105$ with denominators $M^2$ and $2M^2$ in the project notes.
\end{lemma}
\begin{proof}
By CRT, $\rho$ is the product of the local counts at the prime powers $\ell^a\Vert M$. Fix one,
and absorb the $\ell$-adic unit $(M/\ell^a)^{-2}$ into $A$ and $B$ (this does not change
valuations or the count). For a unit $r$ and $x=r+\ell^ay$,
\[
  \frac{Bx+AN/x}{\ell^{2a}}=\frac{Br+AN/r}{\ell^{2a}}+y\,\frac{B-ANr^{-2}}{\ell^{a}}
  +\sum_{n\ge2}(-1)^nANr^{-1-n}\ell^{a(n-2)}y^n .
\]
The last sum is $\Z_\ell$-valued on $\Z_\ell$. So the expression is integral up to a constant on the
class iff $\ell^a\mid Br^2-AN$. Let $t=\min(v_\ell(A),v_\ell(B))$. If $t\ge a$, every unit class
qualifies, and there are at most $\ell^a=\gcd(A,B,\ell^a)$ of them. If $t<a$ and
$v_\ell(A)\ne v_\ell(B)$, then $v_\ell(Br^2-AN)=t<a$ (as $r,N$ are units), so no class qualifies. If
$v_\ell(A)=v_\ell(B)=t<a$, the condition is $r^2\equiv u\pmod{\ell^{a-t}}$ for a unit $u$, which has
at most $2$ solutions modulo $\ell^{a-t}$ for odd $\ell$ (at most $4$ for $\ell=2$), hence at most
$2\ell^t$ (resp.\ $4\cdot2^t$) classes modulo $\ell^a$, and $\ell^t=\gcd(A,B,\ell^a)$. Multiply over
$\ell\mid M$.
\end{proof}

\begin{theorem}[Theorem B; proved in the model, with Lemma~\ref{lem:classcount}]\label{thm:B}
If the degree-one two-term targets of a covering family use $J$ distinct moduli, then
$X\ge cN^{1/5}(JN^{o(1)})^{-1/5}$. So exponent $1/5-\eta_0$ needs $J\ge N^{5\eta_0-o(1)}$ moduli.
\end{theorem}
\begin{proof}
Targets with $M>R$ cover at most one candidate per window, at most $X$ in total, so assume
$M\le R$. For one modulus $M$, by Lemma~\ref{lem:classcount} the number of pairs
(unit class, $(A,B)$) with height $\eta\le h$ is at most
\[
  4^{\omega(M)}\!\!\sum_{\substack{0<|A|,|B|\le h}}\!\!\gcd(A,B,M)
  =4^{\omega(M)}\sum_{e\mid M}\varphi(e)\,\#\{0<|A|,|B|\le h:\ e\mid A,\ e\mid B\}
  \le4^{\omega(M)}\cdot4h^2\sum_{e\mid M}\frac{\varphi(e)}{e^2},
\]
using $\gcd(A,B,M)=\sum_{e\mid\gcd(A,B,M)}\varphi(e)$. The last sum is at most $\tau(M)$, so the count is
at most $4^{\omega(M)}\tau(M)\cdot4h^2=h^2M^{o(1)}=h^2N^{o(1)}$. Over $J$ moduli it is at most $J'h^2$
with $J'=JN^{o(1)}$. Hence for the at most $X$ distinct triples, $\sum_\tau1/\eta_\tau$ is largest
when the heights are the smallest available, which gives $\sum_\tau1/\eta_\tau\le C(J'X)^{1/2}$. By
Proposition~\ref{prop:pertarget}(i) and Cauchy--Schwarz these targets cover at most
$CR^{1/2}(\sum W)^{1/2}(\sum1/\eta)^{1/2}+CX\le CR^{1/2}X(J'X)^{1/4}+CX$ candidates. The other
targets are handled as in Theorem~\ref{thm:A} (if $X\le\delta R^{2/5}$ they cover $o(R)$ or a small
proportion). So $X^{5/4}\ge cR^{1/2}J'^{-1/4}$, i.e.\ $X\ge cR^{2/5}J'^{-1/5}$.
\end{proof}

\begin{theorem}[Theorem C: CRT is multiplicative; proved]\label{thm:C}
Fix a support $I$. For $M\ge1$ with $\gcd(M,N)=1$ and a unit class $r\bmod M$, let $L_{r,M}$ be the
set of coefficient vectors $\alpha\in\Z[1/M]^I$ for which some constant makes the exponent map
satisfy the integrality condition on the class $r$, and let $\Lambda_I=\Z^I$. Let $M=M_1M_2$ with
$\gcd(M_1,M_2)=1$ and $r_j=r\bmod M_j$. Then $\Lambda_I\subseteq L_{r_j,M_j}$,
$L_{r,M}=L_{r_1,M_1}+L_{r_2,M_2}$, $L_{r_1,M_1}\cap L_{r_2,M_2}=\Lambda_I$, and hence
$[L_{r,M}:\Lambda_I]=[L_{r_1,M_1}:\Lambda_I]\,[L_{r_2,M_2}:\Lambda_I]$ (all indices are finite by
Lemma~\ref{lem:K}).
\end{theorem}
\begin{proof}
Integer coefficient vectors satisfy the condition with constant $0$, because $x$ and $N$ are
$\ell$-adic units on unit classes; so $\Lambda_I\subseteq L_{r,M}$ for every $(r,M)$. The integrality
condition is a conjunction of conditions at the primes $\ell\mid M$, and at $\ell\mid M_1$ it depends on
$r$ only through $r\bmod\ell^{v_\ell(M)}=r_1\bmod\ell^{v_\ell(M_1)}$. By partial fractions,
$\Z[1/M]^I/\Z^I=\Z[1/M_1]^I/\Z^I\oplus\Z[1/M_2]^I/\Z^I$, so every $\alpha\in\Z[1/M]^I$ is
$\alpha_1+\alpha_2$ with $\alpha_j\in\Z[1/M_j]^I$, unique modulo $\Z^I$; similarly for the constant. At a
prime $\ell\mid M_1$, the vector $\alpha_2$ is $\ell$-integral and its map is $\ell$-integral on unit
classes, so the condition for $\alpha$ at $\ell$ is equivalent to that for $\alpha_1$; symmetrically at
$\ell\mid M_2$. Hence $\alpha\in L_{r,M}$ iff $\alpha_1\in L_{r_1,M_1}$ and $\alpha_2\in L_{r_2,M_2}$, which gives the sum
decomposition. A vector in both $L_{r_1,M_1}$ and $L_{r_2,M_2}$ has denominators dividing a power
of $M_1$ and a power of $M_2$, hence lies in $\Z^I$. So
$L_{r,M}/\Lambda_I\cong L_{r_1,M_1}/\Lambda_I\oplus L_{r_2,M_2}/\Lambda_I$.
\end{proof}

\begin{remark}[Interpretation]
Theorem~\ref{thm:C} says that a scheme covering \emph{all} classes of a composite modulus gains
nothing from its CRT structure (heuristically, via a volume argument; the rigorous
single-modulus statement is the $J=1$ case of Theorem~\ref{thm:B}). What escapes the $N^{1/5}$
bound is \emph{selective} covering: each $p$ needs only one modulus whose class is reached by
a short weight vector. $\DCL$ uses only the at most two classes per $(v,\ell)$ that have a
root, a $2/\ell$ fraction, but each is reached by a vector of height at most $H$ instead of the
typical $\sqrt\ell$. Across $N^{1/6+o(1)}$ coprime moduli these classes contain the true factor
for almost every pair in $\Sn$ (Theorem~\ref{thm:main}). Whether they cover almost every
candidate for a fixed $N$ is not proved, and covering every candidate is the content of
Conjecture~\ref{conj:DCH}.
\end{remark}

\begin{remark}[Quantifiers]
Theorems~\ref{thm:A}--\ref{thm:B} bound families that cover \emph{every} integer candidate for
a fixed $N$. Theorem~\ref{thm:main} bounds the failure at the \emph{true} factor, averaged over
$\Sn$. So Theorem~\ref{thm:main} does not show that the model exponent is $1/6$ for almost every
$N$. Conversely, Theorem~\ref{thm:A} does not show that an average-case algorithm needs cost
$N^{1/6}$, even within the same sumset framework: an algorithm that only has to reach the true
factor of most inputs is not required to cover every candidate. Whether such algorithms can
go below $1/6$ is open (Section~\ref{sec:open}). Similarly, removing factors up to $N^{1/3}$ by
Pollard--Strassen costs $N^{1/6+o(1)}$; that is the cost of a known method, not a lower bound.
\end{remark}

%======================================================================
\section{Experiments}\label{sec:experiments}
%======================================================================

All experiments were pre-registered: predictions, decision rules, designs and metrics were
committed before any pre-registered run (\texttt{notes/dcltheorem\_prereg.md}, commit
\texttt{05d2bf5}; this was \texttt{3550e24} before a history rewrite, see
\texttt{COMMIT\_SHA\_MAP.md}). The algorithm sees only $N$; the harness samples pairs
uniformly from $\Sn$ and assigns outcome classes after the fact.

\subsection{Designs and metric}
\begin{itemize}
\item \textbf{C1 (primary)} is the algorithm of Section~\ref{sec:alg} (four bands, $\ell\equiv3\bmod4$,
  base $2$, $n$-only parameters) with a finite-size height schedule: $H$ is the smallest
  value with block-0 multiplicity $\lambda_0=|V|\sum_{\ell\in\mathcal P_0}1/\ell\ge3$, which corresponds
  to $H\asymp\sqrt n$. This lies below the range of Corollary~\ref{cor:polylog}; at $H\asymp\sqrt n$
  the proved bound $n^2/H^2$ on the exceptional fraction exceeds $1$ and is vacuous. The
  experiments therefore test correctness, the lemmas and the cost exponent, not the proved
  exceptional rate.
\item \textbf{C2}: C1 with $\lambda_0\ge\ln(P/\ln P)+2$ (heuristic full coverage), $n\le64$.
\item \textbf{C3}: C1 with all odd band primes, Tonelli--Shanks roots and a charged
  least-nonresidue search.
\item \textbf{HH}: a calibrated operational Harvey--Hittmeir implementation, with Bluestein
  chirp powers charged.
\item \textbf{Metric} $E$: $n$-bit modular-multiplication equivalents,
  $E=\mathrm{mul}_N+\mathrm{mul}_{\mathrm{small}}+\mathrm{poly}+\lceil\log_2n\rceil(\mathrm{gcd}_N+\mathrm{inv}_N)
  +\lceil\mathrm{bigmul\ bits}/n\rceil$, with polynomial work charged as
  $\sum(D+1)\lceil\log_2(D+2)\rceil$ (FFT model). $E$ is a proxy; it is not the proof of
  Theorem~\ref{thm:main}(ii). $E_{\mathrm{first}}$ counts the blocks up to the first that yields
  a factor, $E_{\mathrm{full}}$ all four blocks.
\item \textbf{Empirical certificate.} The harness records, for each pair, whether \emph{some}
  block that is clean for $q$ contains \emph{some} target covering $p$ with
  $d\nmid(A-B)(p-1)-\ell^2\delta$ for all $|\delta|<m$, where $d=\ord_q(2)$ and $m$ is the baby length
  actually used.
  This is sufficient for success: such a target's true giant is not FULL by
  Lemma~\ref{lem:full} (the exact $\delta$-range lies inside $|\delta|<m$), so its gcd is $p$ by
  Lemma~\ref{lem:extract}. It is \emph{not} the certificate of Proposition~\ref{prop:cert}, which
  uses the least clean block, the larger range $|\delta|<\mbar$, and requires $\ord_q(2)>T$.
  Below, ``certified'' always means the empirical certificate. (The GPT-6.1-sol round-14
  referee reports that restricting the recorded predicate to the least clean block gives
  $95.8\%$, $97.8\%$, $96.7\%$ and $97.3\%$ in the four certificate datasets; we have not
  re-computed these.)
\end{itemize}

\subsection{Lemma checks (28--40 bits)}
For $n\in\{28,32,36,40\}$, 60 pairs each (design C1, all four blocks recorded): 240/240
correct and factored; 0 violations of Lemmas~\ref{lem:lock}--\ref{lem:modq} on every
covering target; 0 violations of the equivalence ``gcd $=p$ iff not FULL'' (with the exact
$\delta$-range) on 3002 clean covering targets; empirically certified $\Rightarrow$ success in
240/240.

\subsection{Cost (40--80 bits)}
C1 and C3: 32 pairs per size; C2: 8 pairs per size up to 64 bits; HH: 96 runs, 17 excluded
from slopes as pre-registered (16 exact-match early exits, 1 baby-order exit).

\begin{center}
\footnotesize
\begin{tabular}{rrrrrrr}
\toprule
$n$ & $H/Y$ & succ. & C1 $E_{\mathrm{first}}$ & C1 $E_{\mathrm{full}}$ & HH $E$ & C1/HH\\
\midrule
40 & 9/81 & 32/32 & $1.03\cdot10^5$ & $1.07\cdot10^6$ & $2.12\cdot10^4$ & 4.87\\
48 & 9/206 & 32/32 & $3.26\cdot10^5$ & $2.91\cdot10^6$ & $7.57\cdot10^4$ & 4.31\\
56 & 10/503 & 32/32 & $1.01\cdot10^6$ & $9.01\cdot10^6$ & $2.96\cdot10^5$ & 3.40\\
64 & 11/1461 & 32/32 & $4.08\cdot10^6$ & $3.75\cdot10^7$ & $9.43\cdot10^5$ & 4.33\\
72 & 11/3097 & 32/32 & $1.06\cdot10^7$ & $8.84\cdot10^7$ & $3.43\cdot10^6$ & 3.10\\
80 & 11/7804 & 32/32 & $2.76\cdot10^7$ & $2.36\cdot10^8$ & $1.18\cdot10^7$ & 2.35\\
\bottomrule
\end{tabular}
\end{center}

Least-squares slopes of $\log_2(\text{median})$ against $n$ over 48--80 bits:

Columns: design C1 height and band scale, C1 success, medians of $E$, and the ratio of
median $E_{\mathrm{first}}$ (C1) to median $E$ (HH).

\begin{center}
\small
\begin{tabular}{lrrrr}
\toprule
quantity & C1 & C3 & C2 (48--64) & HH\\
\midrule
$E_{\mathrm{first}}$ & \textbf{0.203} & 0.196 & 0.227 & \textbf{0.226}\\
$E_{\mathrm{full}}$ (4 blocks) & 0.200 & 0.202 & 0.237 & --\\
list elements (block 0 / HH list) & 0.171 & 0.163 & 0.189 & 0.196\\
\bottomrule
\end{tabular}
\end{center}

\paragraph{Pre-registered decisions.}
\begin{itemize}
\item \emph{R-a (correctness)} holds: 0 incorrect outputs in 656 DCL runs (the 416 runs of this
  subsection, C1, C3 and C2, plus the 240 lemma-check runs), 0 lemma violations, and every
  empirically certified instance factored.
\item \emph{R-b (exponent)}: $\text{slope}_{C1}-\text{slope}_{HH}=-0.023\le-0.01$, so ``DCL-det has
  a smaller measured exponent'' at these sizes. Both slopes fell inside their predicted
  ranges ($[0.16,0.21]$ and $[0.18,0.23]$).
\item \emph{R-c (constant)} holds as predicted: HH is cheaper at every size, by a factor
  $2.35$--$4.87$. A naive extrapolation of the slope gap puts a crossover near 130 bits. We
  do not claim it: the slope difference partly reflects the $H$-schedule, and the measured
  $E$ slopes exceed the list-element slopes by about $0.03$ for both algorithms
  (polylogarithmic drift).
\item \emph{R-d}: with heuristically full coverage (C2) there is no advantage: slope $0.227$,
  within $0.005$ of HH's (the pre-registered threshold).
\item \emph{R-e}: all three frozen predictions on the certificate components held (below).
\end{itemize}
The four-band worst-case guarantee costs a constant: $E_{\mathrm{full}}/E_{\mathrm{first}}\approx8.5$ at these
sizes. Some first hits are accidental gcd coincidences (83 of 768 C1 blocks, below the null
expectation of 164); any proper divisor is a valid output.

\subsection{Certificate components (48--96 bits)}
For 1000 pairs per size (C1, no collision stage):

\begin{center}
\small
\begin{tabular}{rrrr}
\toprule
$n$ & emp.\ certified & block 0 clean (predicted) & block 0 uncovered ($e^{-\lambda_0}$)\\
\midrule
48 & 1.000 & 0.939 (0.944) & 0.040 (0.049)\\
64 & 1.000 & 0.956 (0.954) & 0.021 (0.024)\\
80 & 1.000 & 0.976 (0.963) & 0.033 (0.046)\\
96 & 1.000 & 0.969 (0.969) & 0.028 (0.029)\\
\bottomrule
\end{tabular}

\medskip
\begin{tabular}{rrrr}
\toprule
$n$ & all 4 blocks uncovered & FULL / covering targets & median bits of $\ord_q(2)$ / of $q$\\
\midrule
48 & 0 & 5 / 11130 & 24 / 25\\
64 & 0 & 0 / 13321 & 32 / 33\\
80 & 0 & 0 / 11273 & 40 / 41\\
96 & 0 & 0 / $\approx1.1\cdot10^4$ & 48 / 49\\
\bottomrule
\end{tabular}
\end{center}

The predicted clean fraction is $1-\sum_{\ell\in\mathcal P_0}1/(\ell-1)$. Uncovered fractions sit
at or below the Poisson value, and the three exceptional sets of the proof are empirically
negligible at these sizes. These data are consistent with the heuristic $e^{-\lambda}$ regime; they
do not test the proved rate.

\subsection{Earlier model-count experiment}
A preceding pre-registered experiment (\texttt{notes/multimod\_results.md} \S6) measured list
elements $m+\sum\lceil W/m\rceil$ for the coverage scheme alone, with a greedy design at about
95\% success: slope $0.171$ over 56--96 bits against $0.196$ for a calibrated HH model count,
and $0.45\times$ the HH count at 96 bits. At matched 100\% success its slope was $0.196$, equal
to HH. Its verdict (SUPPORTED) required a disclosed post-hoc fix of a result classifier.

%======================================================================
\section{Limitations and open problems}\label{sec:open}
%======================================================================

\paragraph{Worst case.} Theorem~\ref{thm:main} is not a worst-case result. Three independent
obstacles remain.
\begin{enumerate}[label=(\alph*)]
\item \emph{Coverage of every candidate by the algorithm's targets.} With polylogarithmic
parameters this is the following conjecture, stated to match the target family of
Section~\ref{sec:alg}.
\begin{conjecture}[Divisor-cover hypothesis, algorithmic form]\label{conj:DCH}
There is $\kappa>0$ such that, with $H=\lceil(\log N)^\kappa\rceil$, $V$ and the bands $\Pk$ as in
Section~\ref{sec:alg}, for all large $N$ coprime to $6$ and every integer
$p\in[\sqrt{N/2},\sqrt N]$ there are $(A,B)\in V$ (so $1\le A,B\le H$, $\gcd(A,B)=1$, $A\ne B$), a
block $k\in\{0,1,2,3\}$ and a prime $\ell\in\Pk$ with $\ell\mid Bp^2-AN$.
\end{conjecture}
The project notes state a weaker-looking variant, DCH$(\kappa)$, which allows $A=B$, all
primes, and the band $(N^{1/6},N^{1/6}(\log N)^\kappa]$; that variant would support a modified
algorithm with those bands and weights, not $\DCL$ as specified. Only the notes' variant was
tested: it is \emph{experimentally supported} at 32--48 bits for greedy designs (zero
uncovered candidates among $1.6\cdot10^4$--$4.9\cdot10^6$ per instance, 25 instances, at
$\lambda\approx\ln\sqrt N+2$). The large sieve gives only ``almost all''. Even full coverage
gives only coverage; extraction needs (b) and (c) as well.
\item \emph{Primes $q$ with small $\ord_q(2)$} (at most $N^{1/4-\varepsilon}$). The proof discards them.
A worst-case version needs a base of large order modulo the specific $q$, on a clean band;
large order modulo $N$ \cite{hh2026} does not give that.
\item \emph{Pairs whose every covering true giant is FULL.} They are rare on average but not
excluded for a fixed $N$.
\end{enumerate}

\paragraph{Unbalanced ratios.} Only $q/p\le2$ is covered. For $q/p\approx\varrho$ the weights
need $B/A\approx\varrho$ and the window grows to about $\sqrt{\varrho N}$; with a shared list a band
costs about $\varrho^{1/4}N^{1/6}$, and summing dyadic ratios up to $N^{1/3}$ gives worse than
Harvey's $1/5$ (heuristic; \texttt{notes/multimod\_results.md} \S3.9). A ratio-uniform version
is open. That the method extends to $q/p\le N^{o(1)}$ is a heuristic expectation, not done
here.

\paragraph{Exceptional rate.} The binding step is the large sieve's $n^2/H^2$. The Poisson
heuristic predicts $e^{-cH^2/n}$, which would give $N^{-\varepsilon}$ exceptions at polylogarithmic
cost (\emph{conjectured}). For a fixed $q$ the $p$-interval has length $P\approx Y^3$, so a sieve
whose level is limited to about $P$ sees only moduli built from at most three band primes.
We \emph{expect} (heuristic, not a theorem) that one-variable upper-bound sieves at such levels
give only polynomial savings in $1/\lambda$. We have not proved any barrier: averaging over $q$
gives a two-variable problem over the box with level up to $P^2$, where bilinear or
dispersion methods might do better.

\paragraph{Below $1/6$.} For \emph{full-candidate covering} in the one-base sumset model with
polynomial exponent maps, Theorem~\ref{thm:A} rules out cost below $N^{1/6}$, whatever the
moduli. Within that model, escaping would need a different collision mechanism
(multi-base searches with non-split known relations are not covered by the reduction in the
project notes). For average-case algorithms, which only need to reach the true factor of
most inputs, we know no such lower bound: whether one can go below $1/6$ on almost all of
$\Sn$ is open. A full factoring algorithm on all inputs would also need small-factor removal
cheaper than Pollard--Strassen's $N^{1/6+o(1)}$ for factors up to $N^{1/3}$.

\paragraph{Practice.} At 40--80 bits the algorithm is $2.3$--$4.9$ times more expensive than
Harvey--Hittmeir in our cost proxy. The asymptotic advantage, if real at larger sizes, is
in the exponent only. The heuristically subexponential methods (ECM, quadratic and number
field sieves) are of course far faster; this work concerns only rigorous deterministic
bounds.

%======================================================================
\section*{Acknowledgements}
%======================================================================

This work was carried out with heavy AI assistance, and we describe the process so that
readers can weigh it. A human ([Author]) set the research question and direction, proposed
several of the ideas that were tested (recorded in the repository under
\texttt{notes/user\_ideas/}), and decided which threads to pursue. An orchestrating Claude
session (Anthropic, Claude Opus 5.5) planned the work and reviewed results; all code,
proofs, experiments and notes were produced by Claude Opus 5.5 subagents working under
written rules (exact arithmetic, no access to the factors in operational code, status labels
on every claim, pre-registration before experiments). This paper was drafted by such a
subagent from the project notes.

GPT-6.1-sol (OpenAI, run through the Codex CLI with web search) was used for literature
sweeps, blind attempts and adversarial review. The following reviews bear directly on this
paper.
\begin{itemize}
\item The round-12 GPT-6.1-sol audit of the predecessor result (an average-case coverage
  theorem in the search model) found a substantive gap (multiplicative orders of the shared
  base $\beta^L$, with the explicit failure example $\beta=3$, $N=1333=31\cdot43$,
  $L=(3\cdot5\cdot7)^2$, where $\beta^L\equiv-1$), a flawed exceptional-set normalisation, and
  missing cost terms. Theorem~\ref{thm:main} was built to close these, and its normalisation
  follows the audit's repair.
\item A GPT-6.1-sol referee report on Theorem~\ref{thm:main} and an independent read-only
  Claude Opus critic, who re-implemented the algorithm from its specification and ran
  adversarial inputs, both returned ``accept with fixes''. The fixes (mandatory baby-length
  rule, missing cost terms, $T(T+1)$ in place of $T^2$, constants, wording) were applied.
\item A fresh round-13 GPT-6.1-sol referee marked every proof step ``correct'' and
  independently re-ran the lemma checks without the repository code. It marked the
  implementation's unrestricted ``never wrong'' claim WRONG (the code returned $N$ itself for
  some prime $N$; fixed and regression-tested) and a heuristic sieve-barrier remark in the
  notes a GAP (now labelled heuristic, Section~\ref{sec:open}). It also located McKee--Pinch
  \S5 and Carella's preprint and asked for the comparison in Section~\ref{sec:prior}; we
  mention its objection to \cite{carella} there without adopting it.
\item This draft itself was reviewed by an independent read-only Claude Opus critic
  (``accept with fixes'') and by a round-14 GPT-6.1-sol journal-style referee (``needs fixes
  first''). Their points were addressed in this revision. Among them: the quantifier scope of
  the lower bound, the unit-class and window-count hypotheses of the model, the proofs of
  Lemma~\ref{lem:K} and Theorem~\ref{thm:C}, the empirical certificate, the algorithmic form
  of Conjecture~\ref{conj:DCH}, and the bibliography.
\end{itemize}
No human has yet checked the proofs line by line. A Lean~4 formalisation of some lemmas was
started at the time of writing; nothing here relies on it.

%======================================================================
\appendix
\section{Reproducibility}
%======================================================================
Code, raw data and notes are in the project repository. The algorithm is
\texttt{src/dcltheorem/algorithm.py}; parameters \texttt{src/dcltheorem/params.py}; tests
\texttt{tests/test\_dcltheorem.py}. The experiments are reproduced by
\begin{verbatim}
.venv/bin/python -m pytest -q tests/test_dcltheorem.py
PYTHONPATH=. .venv/bin/python -m src.dcltheorem.experiments config
PYTHONPATH=. .venv/bin/python -m src.dcltheorem.experiments lemma
PYTHONPATH=. .venv/bin/python -m src.dcltheorem.experiments hh
PYTHONPATH=. .venv/bin/python -m src.dcltheorem.experiments oper
PYTHONPATH=. .venv/bin/python -m src.dcltheorem.experiments exc
PYTHONPATH=. .venv/bin/python -m src.dcltheorem.experiments summarize
\end{verbatim}
Raw records are in \texttt{experiments/raw/dcltheorem\_*.jsonl} and the summary in
\texttt{experiments/analysis/dcltheorem\_summary.json}. Post-run code changes (an integer-only
$m$-grid, an integer cube root, GMP products in the exponent tree, a finite-size guard, and
the output validation) do not change any recorded run; twelve recorded C1/C3 instances and
four recorded C1 instances were re-run with identical results.

\bibliographystyle{plain}
\bibliography{refs}

\end{document}
